% 张三慧第4版第二卷第29章《固体中的电子》思考题29.1--29.13
% 主教材V2扫描PDF物理页452，须独立计数，不能与同编号习题混淆。
\section*{第29章\quad 固体中的电子：思考题}
\paragraph{思考题 29.1}
金属中的“自由”电子是在什么近似条件下可以看成自由的？晶体的周期势如何影响这一近似？
\paragraph{思考题 29.2}
为什么金属的传导电子对比热的贡献往往较小，却能带来很大的电导率？
\paragraph{思考题 29.3}
量子简并费米气电子在能量空间中的分布由什么原理决定？与经典Maxwell--Boltzmann气体有何不同？
\paragraph{思考题 29.4}
什么是电子能带、禁带、导带、价带？如何区分孤立原子的分立能级与晶体的能带？
\paragraph{思考题 29.5}
金属、绝缘体和半导体在能带和费米能级结构上有何关键差异？
\paragraph{思考题 29.6}
硅晶体以磷取代少量硅原子后产生什么型半导体？
移动电子与空穴浓度的变化如何？电中性条件是否被破坏？
\paragraph{思考题 29.7}
锗晶体中掺入硼后空穴浓度增大，导带中的自由电子浓度是否必然相应增大？
空穴浓度增加会使晶体净带正电吗？
\paragraph{思考题 29.8}
本征半导体、单一杂质半导体和PN结是否都具有同一种单向导电特性？
区分体材料电导与整流结效应。
\paragraph{思考题 29.9}
在霍尔效应测量中，为什么半导体薄片通常比普通金属片更容易产生可测的霍尔电压？
\paragraph{思考题 29.10}
水平矩形N型半导体片长边沿东西向，
约定正电流由北向南，外加磁场竖直向上。
求其霍尔电场在水平面内的方向；
把材料改为P型而维持电流、磁场不变，霍尔电场方向如何变化？
\paragraph{思考题 29.11}
对于电子和空穴共存的本征半导体，能否测量到霍尔电压？
说明两个载流子贡献抵消的条件与不同迁移率的影响。
\paragraph{思考题 29.12}
在图29.17的N型沟道MOSFET中，
保持源漏偏压，逐渐增大源极相对栅极的电压 \(U_{SG}\)，
直到导电通道被耗尽而 \(I_{DS}\to0\)。
通道是先在源端、漏端还是同时处处被阻断？
说明对局部沟道电势分布的依赖。
\paragraph{思考题 29.13}
普通电视机或类似电器的红外遥控系统为何需要半导体元件？
遥控器的发射端和电视机的接收端各可使用什么类型的半导体器件？
